\section{The arithmetic interface and a precise transfer principle}
\label{sec:arithmetic}

\subsection{The object theory}

There is one numerical sort. For every specified primitive recursive
definition there is a function symbol. The formation rules are zero,
successor, projections, composition, and primitive recursion, with equations
\[
 r(0,\boldsymbol x)=g(\boldsymbol x),\qquad
 r(t+1,\boldsymbol x)=h(t,r(t,\boldsymbol x),\boldsymbol x).
\]
The nonlogical axioms include these equations, $Sx\ne0$, successor
injectivity, and
\[
 A(0,\boldsymbol x)\to
 \bigl(\forall t(A(t,\boldsymbol x)\to A(t+1,\boldsymbol x))\bigr)
 \to\forall t\,A(t,\boldsymbol x)
\]
for quantifier-free $A$. Equality has reflexivity and substitution. The
logical rules are intuitionistic first-order natural deduction: introduction
and elimination for $\wedge,\vee,\to,\forall,\exists$, falsity elimination,
and the usual freshness conditions. No excluded-middle schema, choice,
Markov principle, general comprehension, or reflection rule is added.
This is the presentation implemented by \texttt{PR}, \texttt{PTm},
\texttt{PRAAx}, and \texttt{Prf} in the inspected repository
\cite{repo-pr,repo-deduction}. Avigad's lectures give the general
proof-theoretic background \cite{avigad2015}.

The restriction concerns the \emph{induction formula}. Quantified assumptions
and conclusions are allowed. For example, a hypothesis $\forall k\,R(k)$
can be instantiated at the finitely many indices needed in a computation;
one must not then use $\forall k\,R(t,k)$ as an induction invariant unless a
separate admissibility argument has been supplied.

\subsection{Coding and totalization}

Fix explicit numerical encodings. One convenient list encoding is
\[
 []=0,\qquad \operatorname{cons}(a,t)=2^a(2t+1).
\]
For a positive code, bounded divisibility tests recover its exponent of two
and then its tail. The tail is strictly smaller than the code. A loop bounded
by the original code therefore recovers length and entries. Pairing, finite
tables, tagged unions, and arrays are implemented by the same operations.
Efficiency is irrelevant to the existence of a PR definition, although the
reference code uses ordinary lists for practicality.

Integers use a sign and magnitude; rationals use an integer numerator and a
positive denominator. Equality is cross multiplication. Gcd normalization is
optional for basic arithmetic, but is useful for the rational denominator
argument in Section~\ref{sec:macaulay}. A gcd can be obtained by scanning
bounded possible divisors, or by a bounded Euclidean calculation. Division and
remainder by a positive integer are PR; a zero divisor is assigned a default
value. Rational inversion is totalized and its field laws are used only with
the nonzero-numerator guard. Malformed codes also receive default values.
Every theorem retains the guard that makes those values irrelevant.

Uniform evaluation of a finite polynomial table is legitimate: enumerate its
finitely many monomials, compute each finite product of powers, and add its
coefficient. Likewise a finite function table can be looked up uniformly.
Neither operation evaluates arbitrary PR \emph{programs}. The definition of
the evaluator is fixed; only a finite data table is variable.

\subsection{Booleanization and bounded choice}

\begin{lemma}[Bounded Booleanization]
\label{lem:booleanization}
For each fixed bounded numerical formula $B(\boldsymbol x)$ whose atoms compare
fixed PR terms, there is a PR Boolean $b_B$ and an iPRA derivation
\[
 B(\boldsymbol x)\ \longleftrightarrow\ b_B(\boldsymbol x)=1.
\]
Hence bounded induction is an admissible use of quantifier-free induction.
The statement is a schema over fixed syntax, not a universal truth evaluator.
\end{lemma}

\begin{proof}
Natural-number equality and order have characteristic functions whose
specifications follow from the successor axioms and elementary QF inductions.
Signed and rational comparisons reduce to them. Boolean conjunction,
disjunction, and complement handle the logical operations. For implication,
use $\neg b\vee c$ and case analysis on the already decidable Boolean $b$.
This does not introduce excluded middle for an arbitrary formula.

For bounded universals define a prefix conjunction
$p(0)=1$, $p(t+1)=p(t)b_C(\boldsymbol x,t)$. Fix a terminal bound $T$.
Under $\forall j<T\,C(\boldsymbol x,j)$, induct on the QF invariant
$t\le T\to p(t)=1$, using the hypothesis only at $j=t$.
Conversely, for a fixed numerical parameter $j$, induct on
$j<t\wedge p(t)=1\to b_C(\boldsymbol x,j)=1$.
The successor case is the decidable split $j=t$ or $j<t$.
After the induction, generalize $j<T$. This proves the bounded-universal
equivalence without assuming that equivalence as part of the invariant.

For bounded existence, maintain a prefix disjunction $e(t)$ and a first-success
register $w(t)$. If no success has been seen, the register is the current
sentinel $t$; otherwise it retains the first successful index. On appending a
new test at $t$, retain an earlier success, record $t$ if the new test succeeds,
and otherwise change the sentinel to $t+1$. The QF assertions
\[
 e(t)=1\to w(t)<t\wedge b_C(\boldsymbol x,w(t))=1,
 \qquad j<t\wedge b_C(\boldsymbol x,j)=1\to e(t)=1
\]
follow by induction, with $j$ fixed in the second assertion. They prove both
directions and supply an actual bounded witness. Structural recursion on the
fixed formula completes Booleanization. Replace a bounded induction formula by
its Boolean equation, apply QF induction, and translate back.
\end{proof}

In particular, if $R$ is a bounded test and $b(x,i)$ a supplied PR bound, then
\[
 \forall i<n\ \exists y<b(x,i)\,R(x,i,y)
\]
admits a PR finite list of witnesses, obtained by bounded least-witness search.
There is no corresponding inference from unbounded existence without a further
constructor or bound. This difference will reappear in equivalence of
categories and finite skeleton extraction.

\subsection{Proofs of bounded loops}

\begin{lemma}[Guarded loop rule]
\label{lem:loop}
Let $S(x,0)=a(x)$ and $S(x,t+1)=s(x,t,S(x,t))$ be fixed PR definitions.
Let $I(x,N,t,z)$ and $Q(x,N,z)$ be bounded predicates. Suppose iPRA already
derives, from a context $\Gamma$ independent of the step-local variables
$t,z$, the initialization,
transition, and exit assertions
\begin{align*}
 I(x,N,0,a(x)),\qquad
 t<N\wedge I(x,N,t,z)&\to I(x,N,t+1,s(x,t,z)),\\
 I(x,N,N,z)&\to Q(x,N,z).
\end{align*}
Then iPRA derives $Q(x,N,S(x,N))$ from $\Gamma$.
\end{lemma}

\begin{proof}
Under the fixed context, induct on the Booleanization of
$t\le N\to I(x,N,t,S(x,t))$. If $t+1\le N$, then $t<N$ and the transition
applies. If the guard fails, the implication holds by decidable arithmetic.
At $t=N$ the exit assertion gives the result. An unbounded precondition may
remain in $\Gamma$; it must not be silently included in the QF induction
formula. All conditions on changing intermediate data belong in $I$ as
bounded numerical tests.
\end{proof}

Every main construction below comes with such a finite invariant: digit
prefixes for tables, a filtered prefix for finite limits, bounded reachability
for finite colimits, matrix identities and pivot patterns for linear algebra,
and a checked prefix for final enumeration. The theorem gives a way to
\emph{compose existing local proofs}. It does not replace their mathematical
content.

\subsection{Finite diagram proof transport}

\begin{theorem}[Syntactic transport of local mathematical arguments]
\label{thm:transport}
Fix numerical presentations of finitely many operations and explicit iPRA
derivations of the local laws used in a finite proof. Suppose that proof uses
only typed substitution, equality reasoning, intuitionistic first-order rules,
and bounded loops certified as in Lemma~\ref{lem:loop}. Then its translated
conclusion has an iPRA derivation from the translated hypotheses. For purely
equational proof trees, the proof transformation on finite numerical proof
codes is PR when the syntax operations and supplied proof constructors are PR.
\end{theorem}

\begin{proof}
Translate each operation to its specified numerical symbol. Translation
commutes with substitution, by recursion on the finite term. Process the
proof tree from its leaves. Reflexivity, symmetry, transitivity, and congruence
become the corresponding equality rules. A law instance is its supplied
derivation after simultaneous substitution and weakening. If it has a typing
guard, first translate the proof of that guard and then use implication
elimination. Logical rules translate to the same intuitionistic rules, keeping
their freshness conditions. A bounded-loop node is replaced by the QF induction
constructed in Lemma~\ref{lem:loop}. Each node therefore has a derivation once
its predecessors do.

For the stated numerical transformation, store a proof as a list of nodes
whose dependencies have smaller indices. A primitive recursion through this
list builds the translated node using the already built predecessors.
The restrictions on syntax and local proof constructors make each update PR.
No evaluation of an arbitrary program or assertion that every checked proof is
true is used.

If mathematical equality is a relation $E(a,b)$ rather than literal equality
of codes, use the supplied derivations of reflexivity, symmetry, transitivity,
and operation congruence for $E$. Arithmetic Leibniz substitution alone does
not justify replacing one representative by another. This is why the setoid
laws are part of the presentation.
\end{proof}

An immediate practical benefit is that the proof of inverse uniqueness needs
only five local equations, regardless of the category's size. The determining
calculation in Yoneda needs one naturality instance and one unit law. The
accompanying Lean module implements exactly these local proof patterns at the
\texttt{Prf PRAAx} interface; the arithmetic validation of a concrete category
is a separate input.

\subsection{The roles of categorical language}

In the rest of the paper there are three precise levels. A \emph{uniform finite
theorem} quantifies over numerical tables and uses fixed PR table algorithms.
A \emph{presentation-wise theorem} is instantiated after its finitely many PR
symbols have been fixed. An \emph{external category} organizes these instances
in the metatheory. A theorem about that external organization is not, merely
by its wording, a new object or quantification principle inside first-order
iPRA. Keeping these levels explicit allows useful category theory without
silently strengthening the arithmetic foundation.
